\secnum{7}
\section{Comparator with Hysteresis (Schmitt Trigger)}\label{sec:schmitt}
\given{fig5_14}{0.45\linewidth}{Schmitt trigger in PSPICE\label{fig:5-14}}

\subsection{Tính toán lý thuyết}
Hồi tiếp dương nên ngõ ra bão hòa: $V_{OUT} = \pm V_{sat} = \pm 5\,\text{V}$. Ngõ $(+)$ là cầu chia áp từ ngõ ra:
\[
V_{(+)} = V_{OUT}\,\frac{R_{23}}{R_{22} + R_{23}} = 0.8\,V_{OUT} = \pm 4\,\text{V}.
\]
Ngõ ra đổi trạng thái khi $V_{IN}$ vượt $V_{(+)}$:
\[
V_{TH} = +4\,\text{V}\;(V_{OUT} = +5\,\text{V}),\qquad
V_{TL} = -4\,\text{V}\;(V_{OUT} = -5\,\text{V}),\qquad
V_H = V_{TH} - V_{TL} = 8\,\text{V}.
\]
Thời điểm chuyển ($T = 20\,\text{ms}$): $5\sin(2\pi\cdot 50\,t) = 4$
\[
t_1 = \frac{\arcsin 0.8}{2\pi \cdot 50} \approx 2.95\,\text{ms}\;(\text{xuống}),\qquad
t_2 = t_1 + \frac{T}{2} \approx 12.95\,\text{ms}\;(\text{lên}).
\]

\begin{dapsobox}
$V_{OUT} = \pm 5\,\text{V}$, \quad $V_{TH} = +4\,\text{V}$, \quad $V_{TL} = -4\,\text{V}$, \quad $V_H = 8\,\text{V}$.
\end{dapsobox}

\subsection{Mô phỏng PSpice}
\simcfg{Time Domain (Transient)}{Run To Time $=100\,\text{ms}$, Maximum Step Size $=0.1\,\text{ms}$}
\simresult
\simfig{ex37_tran}{Kết quả Time Domain: $V_{IN}$, $V_{(+)}$ và $V_{OUT}$ theo thời gian}

\begin{table}[H]
\centering
\begin{tabular}{l c c}
\toprule
\textbf{Đại lượng} & \textbf{Lý thuyết} & \textbf{Mô phỏng}\\
\midrule
$V_{OUT}$ mức cao / thấp (V) & $+5$ / $-5$ & $+5.000$ / $-5.000$\\
$V_{(+)}$ mức cao / thấp (V) & $+4$ / $-4$ & $+4.000$ / $-4.000$\\
Ngưỡng trên $V_{TH}$ (V)     & $+4$        & $\approx +4$ (3.96)\\
Ngưỡng dưới $V_{TL}$ (V)     & $-4$        & $\approx -4$\\
Thời điểm chuyển xuống (ms)  & 2.95        & $\approx 2.93$\\
Thời điểm chuyển lên (ms)    & 12.95       & $\approx 13.0$\\
\bottomrule
\end{tabular}
\caption{Thông số Schmitt trigger}
\end{table}

\subsection{Giải thích hiện tượng}
\begin{remarkbox}[Giải thích tín hiệu ngõ ra]
\begin{itemize}[nosep]
  \item Ngõ ra là sóng vuông: không có hồi tiếp âm nên op-amp bão hòa ở $\pm 5\,\text{V}$ (bằng \texttt{VPOS}/\texttt{VNEG}).
  \item Ngưỡng so sánh là $V_{(+)} = 0.8\,V_{OUT}$, không cố định ở $0$. Khi $V_{OUT} = +5\,\text{V}$ thì $V_{(+)} = +4\,\text{V}$: ngõ ra chỉ lật xuống khi $V_{IN} > +4\,\text{V}$; ngay sau đó $V_{(+)} = -4\,\text{V}$ nên ngõ ra chỉ lật lên khi $V_{IN} < -4\,\text{V}$. Vì vậy ngõ ra chuyển trạng thái ở $+4\,\text{V}$ và $-4\,\text{V}$.
  \item Trong vùng $-4\,\text{V} < V_{IN} < +4\,\text{V}$ ngõ ra giữ trạng thái cũ (hysteresis), giúp chống nhiễu quanh ngưỡng.
  \item Kết quả mô phỏng khớp lý thuyết: $V_{OUT} = \pm 5\,\text{V}$, $V_{(+)} = \pm 4\,\text{V}$; ngõ ra chuyển xuống tại $\approx 2.93\,\text{ms}$ khi $V_{IN} \approx +4\,\text{V}$ (lý thuyết $2.95\,\text{ms}$) và chuyển lên tại $\approx 13\,\text{ms}$ khi $V_{IN} \approx -4\,\text{V}$ (lý thuyết $12.95\,\text{ms}$), lặp lại sau mỗi $20\,\text{ms}$. Trong khoảng $3\,\text{ms}$ đầu, ngõ ra tăng dần từ $0$ (điểm làm việc ban đầu tại $t = 0$, khi $V_{IN} = 0$) lên $+5\,\text{V}$, rồi lật xuống khi $V_{IN}$ vượt $+4\,\text{V}$ lần đầu.
\end{itemize}
\end{remarkbox}
